\documentclass[11pt]{article}
\usepackage{amsmath,amssymb,amsthm}
\usepackage[margin=1.1in]{geometry}
\newtheorem{definition}{Definition}
\newtheorem{conjecture}{Conjecture}
\theoremstyle{remark}
\newtheorem{remark}{Remark}

\title{The observer ladder:\\ stratifying the two-skeleton conjecture, with a second first light}
\author{Carl Gribble\thanks{Programme Note 3, draft v0.1. Companion to the essay \cite{Essay} and
Programme Notes 1--2 \cite{Defs, Engine}. Developed in collaboration with Claude (Anthropic); the
scanner runs are preregistered as described in Section~4. Responsibility for the content rests with
the author. A full disclosure of AI use appears at the end of the note.}}
\date{August 2026 --- draft}

\begin{document}
\maketitle

\begin{abstract}
The essay's two-skeleton conjecture --- the multiplication crystal's detectable order is its
congruence skeleton and its spectral skeleton, and nothing else --- is here stratified by
\emph{observer complexity}, so that its known base camps become theorems of the ladder and its
open floors become the scanner's hunting ground. The base camps are real: automatic sequences
cannot correlate with the M\"obius signal (M\"ullner, after Mauduit--Rivat's digits of primes),
nor can bounded-depth circuits (Green), and the two-point logarithmically-averaged Chowla
correlation vanishes (Tao), with the zero-entropy stratum largely settled in logarithmic density
(Frantzikinakis--Host). Above the theorem line, the skeleton-null scanner of the essay is upgraded
and rerun under a preregistered protocol: a twelve-feature battery tagged by observer class,
nulls stratified by residue mod~$60$ (closing the mod-$4$ blindness the first instrument found in
itself), survival requiring replication in both halves, and --- new --- a skeleton-matched
\emph{control population} $m = p + 30$, identical to the primes in congruence class, size, and gap
chain but overwhelmingly composite, which automates the audit question ``is this reading
prime-specific?''. Second first light, honestly reported: the v2.0 instrument nulled its own
congruence controls by over-stratification (the scanner finding its own bug, again) and was
revised to v2.1 before any prediction was scored; under v2.1 all five committed predictions hold
--- the Lemke Oliver--Soundararajan repulsion fires at $z = -37.8$ mod $4$ and $z = -41.3$ mod $3$
(the second a new positive control), the gap anticorrelation at $z = -10.8$, the automatic-class
theorem-null is clean, and the v1 chi-cluster is absorbed by the finer null as predicted. Of $57$
survivors, $43$ are reproduced by the control (not prime-specific), $11$ are Pratt-height
entanglements invisible to the control by definition, and $3$ small gap-by-smoothness readings
remain in audit. Confirmed candidates: zero. The conjecture, the ladder, and the protocol are
stated so that any future survivor knows exactly which floor it must climb.
\end{abstract}

\section{The conjecture, recalled}

The essay \cite{Essay} states the two-skeleton conjecture: for every finite pre-registered battery
of crystal-native invariants of the primes, every correlation detectable at any scale is derivable
from the congruence data together with the Hardy--Littlewood correlations \cite{HardyLittlewood};
conditioned on these, all remaining statistics agree with the skeleton-null ensemble. Its random
side is the Chowla--Sarnak circle \cite{Chowla, Sarnak}. As stated, the conjecture is operational
--- it names its own experiment --- but it flattens a structure that the known partial results
articulate: orthogonality is \emph{proved} against weak observers and \emph{open} against strong
ones. This note supplies the missing coordinate.

\section{The ladder}

\begin{definition}[Observer class; orthogonality]
An \emph{observer class} $\mathcal{C}$ is a family of bounded sequences $u \colon \mathbb{N} \to
[-1, 1]$ with a specified membership criterion (automaton size, circuit depth, dynamical entropy,
computation time). The class \emph{detects nothing in} a bounded arithmetic signal $f$ if
$\frac1N \sum_{n \le N} f(n)\, u(n) \to 0$ for every $u \in \mathcal{C}$ (in ordinary or
logarithmic density, as stated). The \emph{ladder} is the family of such statements as
$\mathcal{C}$ grows.
\end{definition}

For the M\"obius/Liouville signal the base camps are theorems:

\begin{center}\small
\begin{tabular}{lll}
\hline
observer class & status for $\mu$ (or $\lambda$) & attribution \\
\hline
automatic sequences & orthogonal --- theorem & M\"ullner \cite{Mullner}; digits of primes:
Mauduit--Rivat \cite{MauduitRivat} \\
$\mathrm{AC}^0$ circuits & orthogonal --- theorem & Green \cite{GreenAC0} \\
two-point self-correlations & vanish, log density --- theorem & Tao \cite{TaoChowla} \\
zero-entropy dynamics & log-Sarnak --- theorem (structured cases) & Frantzikinakis--Host
\cite{FH} \\
polynomial-time observers & open & (pseudorandomness frontier) \\
all bounded observers & open --- Sarnak's conjecture & \cite{Sarnak} \\
\hline
\end{tabular}
\end{center}

\begin{conjecture}[Two-skeleton, stratified]
For each observer class $\mathcal{C}$ in the ladder, every $\mathcal{C}$-detectable correlation
among crystal-native invariants of the primes is derivable from the congruence skeleton together
with the Hardy--Littlewood correlations. The base camps above are the proven strata; the
conjecture is the assertion that the pattern continues to the top of the ladder.
\end{conjecture}

The scanner probes \emph{above} the theorem line: its crystal-native features ($\Omega(p \pm 1)$,
Pratt height, largest-factor slope) are not computable by weak observers, so no theorem protects
them; only the skeleton-null ensemble does.

\section{The instrument, upgraded}

Scanner v2 extends the essay's first instrument \cite{Essay} in four ways, all preregistered
before first light (Section~4). \emph{Battery}: twelve features tagged by observer class ---
eight crystal/spacing features from v1, the characters mod $4$ and mod $3$ (congruence class),
and the binary digit sum, Thue--Morse sign, and Rudin--Shapiro sign (automatic class, riding
directly above the M\"ullner--Mauduit--Rivat base camp). \emph{Nulls}: stratification by coprime
residue mod $60$ crossed with four size bands --- mod $60$ determines the residue mod $4$ and mod
$3$, closing the blindness that v1's own audit discovered when its $2$-adic cluster poured
through a mod-$30$ null. \emph{Survival}: full-sample $|z| \ge 3.9$ \emph{and} same-sign
$|z| \ge 3.0$ in both split halves. \emph{Control}: the population $m = p + 30$ carries the
primes' exact mod-$60$ classes, sizes, and gap chain while being overwhelmingly composite; the
whole battery (except Pratt height, undefined off the primes) runs on it, and a survivor
reproduced by the control with the same sign is classified \emph{not prime-specific} ---
the audit ladder's ``skeleton-derivable or mechanical'' stage, automated.

\section{Preregistration and second first light, honestly}

The registry protocol (directory \texttt{registry/} of \cite{Code}) commits battery, nulls,
thresholds, exclusions, and predictions to version control \emph{before} a run; the commit is the
commitment device, the prereg file is never edited, and results are appended separately. Scanner
v2's preregistration (commit \texttt{930a091}) committed seven predictions, among them: the
consecutive-residue repulsion fires mod $4$ (the v1 positive control) \emph{and} mod $3$ (a new
one); no automatic-class feature survives as prime-specific (the theorem-null); the v1 same-prime
chi cluster is absorbed by the finer stratification; zero candidates after audit.

First light v2.0 immediately caught its own instrument in two errors, which the registry
discipline requires reporting rather than dissolving: over-stratification made the congruence
features constant within strata, nulling exactly the positive controls the prereg predicted would
fire (the mod-$60$ upgrade, over-applied, blinded the instrument to the signal it was built to
confirm); and the theorem-null checker enforced a stricter carve-out than the prereg's text. One
genuinely instructive reading appeared: Rudin--Shapiro against itself at consecutive primes,
$z = -12.8$, replicated --- consecutive primes share binary prefixes, so the automatic sequence's
\emph{own} shift-autocorrelation, an integer phenomenon, leaks through prime gaps. The v2.1
revision (two-tier standardization; consecutive nulls permuted within size bands; the control
population) was made before any prediction was scored, and the control settled the
Rudin--Shapiro reading exactly as diagnosed: reproduced off the primes, not prime-specific.

Under v2.1, all committed predictions hold. The Lemke Oliver--Soundararajan repulsion \cite{LOS}
fires at $z = -37.8$ (mod $4$) and $z = -41.3$ (mod $3$); the gap anticorrelation at $z = -10.8$; the automatic class is clean of
prime-specific survivors; the v1 chi cluster is gone. One instructive subtlety: the control
reproduces the repulsion ($z = -38.5$, $-41.5$) --- as it must, since $m = p + 30$ inherits the
primes' residue \emph{chain}, and the repulsion lives in the skeleton-plus-Hardy--Littlewood
layer that the control copies. The control distinguishes prime-specific structure from
skeleton-borne structure; it does not, and should not, erase the skeleton.

Of the $57$ survivors: $43$ control-reproduced (integer smoothness autocorrelations at small
shifts, size-budget entanglements, and the skeleton-borne controls above); $11$ involving Pratt
height, which the control cannot audit --- their same-prime members are definitional
entanglements (Pratt height is built from the factorization of $p - 1$), their consecutive
members plausibly inherit the control-reproduced smoothness autocorrelations through that
entanglement; and $3$ small readings ($|z| = 5.0$--$5.4$) pairing the normalized gap with
smoothness of neighbouring shifted primes, where the control \emph{disagrees} --- genuinely
prime-only observables, retained in audit for the next scale, not promoted. Confirmed candidates:
zero, as predicted; the honest harvest of a second first light, with the instrument now two bugs
wiser.

\section{What would refute what}

A robust, replicated, control-clean, prime-specific correlation of an automatic-class feature
would contradict no theorem directly --- the base camps govern $\mu$-orthogonality and
single-sequence equidistribution, not consecutive-prime pair correlations --- but it would breach
the stratified conjecture at a floor adjacent to proven ground, and would therefore be
extraordinary. A crystal-class survivor clearing the full ladder would refute the two-skeleton
conjecture as stated. Nothing of either kind has appeared. The undecidability of the weld
(essay, Section~5's logic column) guarantees the hunt cannot be exhaustively closed; the ladder
guarantees any find will know its floor.

\section*{Disclosure of AI use}

This note was developed in an extensive collaboration with Claude (Anthropic; Claude Fable 5,
accessed through Claude Code, 2026). The direction of the work and all decisions about scope and
claims are the author's; the AI system drafted the text, wrote and ran the instrument under the
registry protocol, and assembled the attributions. The base-camp theorems are cited, not
reproved; the instrument revisions and their causes are reported as they occurred. The author has
reviewed the note and takes full responsibility for its content.

\begin{thebibliography}{99}
\bibitem{Essay} C. Gribble, The multiplication crystal: dimension, diffraction, and the
symmetries of arithmetic, manuscript (2026), in the code repository \cite{Code}.
\bibitem{Defs} C. Gribble, Crystals and their shadows, programme note (2026), in \cite{Code}.
\bibitem{Engine} C. Gribble, The multiplication crystal's engine is a formal logarithm, programme
note (2026), in \cite{Code}.
\bibitem{HardyLittlewood} G.~H. Hardy and J.~E. Littlewood, Some problems of `Partitio numerorum'
III, \emph{Acta Math.} \textbf{44} (1923), 1--70.
\bibitem{Chowla} S. Chowla, \emph{The Riemann Hypothesis and Hilbert's Tenth Problem}, Gordon and
Breach, New York, 1965.
\bibitem{Sarnak} P. Sarnak, Three lectures on the M\"obius function, randomness and dynamics,
lecture notes, Institute for Advanced Study, 2011.
\bibitem{MauduitRivat} C. Mauduit and J. Rivat, Sur un probl\`eme de Gelfond: la somme des
chiffres des nombres premiers, \emph{Ann. of Math.} \textbf{171} (2010), 1591--1646.
\bibitem{Mullner} C. M\"ullner, Automatic sequences fulfill the Sarnak conjecture, \emph{Duke
Math. J.} \textbf{166} (2017), 3219--3290.
\bibitem{GreenAC0} B. Green, On (not) computing the M\"obius function using bounded depth
circuits, \emph{Combin. Probab. Comput.} \textbf{21} (2012), 942--951.
\bibitem{TaoChowla} T. Tao, The logarithmically averaged Chowla and Elliott conjectures for
two-point correlations, \emph{Forum Math. Pi} \textbf{4} (2016), e8.
\bibitem{FH} N. Frantzikinakis and B. Host, The logarithmic Sarnak conjecture for ergodic
weights, \emph{Ann. of Math.} \textbf{187} (2018), 869--931.
\bibitem{LOS} R.~J. Lemke Oliver and K. Soundararajan, Unexpected biases in the distribution of
consecutive primes, \emph{Proc. Natl. Acad. Sci. USA} \textbf{113} (2016), E4446--E4454.
\bibitem{Code} C. Gribble, Reference implementations accompanying this programme, code repository
(2026), \texttt{https://github.com/carlgribble-caa/prime-moments}.
\end{thebibliography}

\end{document}
